\begin{abstract}
This paper tests Vibe Theory in simulation and assembles the results into a single picture. The framework holds that reality is one substance, the \emph{vibe}, existing only as a mesh of relations (\emph{notes}) among states (\emph{tones}), with no background space or time, and that everything else is a large-scale pattern of that mesh. We built a finite, seeded, reproducible testbed and posed nine load-bearing questions, then measured the answers. Geometry is recovered sharply and almost uniquely from a discrete causal order, and, with a correct uniform-measure sampler, the substrate makes smooth low-dimensional spacetime a \emph{stable phase} that the action supports, coexisting with the layered phase in the manner of a first-order transition, where without the action the smooth order decays. Whether that phase strictly dominates the sum over histories is left open. Matter is built from the same tones, with spin appearing as a topological invariant of the mesh and a single chiral fermion realized through the overlap construction. Force lives on the links as a gauge connection that confines, carries a topological index equal to the gauge charge, and forms a chiral condensate in dynamical Abelian and non-Abelian fields. The gravitational action is the discrete Benincasa-Dowker functional whose continuum limit is curvature, and it is exactly the term that makes spacetime stable. A hyperbolic random mesh is shown to be Lorentz-safe, navigable without addressing, and stable under growth. Quantum violation is shown to be reachable from the deterministic mesh, with the precise residual tension made quantitative. We give the full epistemic scoreboard, separating what is demonstrated, what has strong evidence but no proof, and what is genuinely open or beyond reach. The conclusion is that monism is computationally coherent and that the dynamics supports a stable spacetime phase, with the strict-dominance, quantum, and chiral-gauge questions stated precisely and left open. This is a simulable model with a concrete build target, not a finished empirical theory.
\end{abstract}
